\documentclass[11pt,letterpaper]{article}
\usepackage[T1]{fontenc}
\usepackage[utf8]{inputenc}
\usepackage{lmodern,microtype,amsmath,amssymb,amsthm,mathtools,booktabs,tabularx,array}
\usepackage[margin=0.89in,top=0.83in,bottom=0.86in,headheight=14pt]{geometry}
\usepackage[dvipsnames]{xcolor}
\usepackage{fancyhdr,enumitem}
\usepackage{hyperref}
\definecolor{ink}{HTML}{163D50}
\definecolor{accent}{HTML}{16635F}
\hypersetup{colorlinks=true,linkcolor=ink,urlcolor=accent,citecolor=accent,
 pdftitle={Three shorter steps in the theta-lift proof of Dirichlet GRH},
 pdfsubject={Auxiliary quadratic twists, a Perron integral, and the dimension of theta-lift images}}
\pagestyle{fancy}\fancyhf{}\fancyhead[L]{\small\color{ink}THETA LIFTS AND DIRICHLET ZEROS}
\fancyhead[R]{\small Proof simplifications}\fancyfoot[C]{\small\thepage}
\setlength{\parindent}{0pt}\setlength{\parskip}{0.48em plus .07em minus .1em}
\setlength{\emergencystretch}{2em}
\setlist{nosep,leftmargin=1.65em}
\newtheorem{theorem}{Theorem}
\newtheorem{lemma}[theorem]{Lemma}
\newcommand{\Source}[3]{\href{https://github.com/VladimirReshetnikov/ProveIt/blob/d8415e76ad850651f8c0083afe689dd9f4a45a10/Analysis/GeneralizedRiemannHypothesis/vendor/grh-lean/grh/src/#1.lean\#L#2}{#3}}
\newcommand{\Half}{\tfrac12}
\newcommand{\E}{\mathcal E}
\newcommand{\sg}{\sigma_0}
\newcommand{\M}{\mathcal M}
\newcommand{\e}{\varepsilon}
\raggedbottom
\begin{document}
\thispagestyle{plain}
{\small\bfseries\color{accent} RESEARCH NOTE\hfill 27 NOVEMBER 2026}\par
\vspace{0.6em}
{\LARGE\bfseries\color{ink} Three shorter steps in the theta-lift proof\\[.15em] of Dirichlet GRH}\par
\vspace{0.3em}
{\large Quadratic twists, a smoothed prime sum, and the dimension of the lifted space}\par
\vspace{0.5em}
\textbf{The proposed changes.}
Alpöge and Claude's argument \cite{source,origin} proves a one-sided estimate for primes inert in an auxiliary imaginary quadratic field. To exclude a proposed zero, it must choose the field so that twisting by its quadratic character leaves the zero detectable. The most economical replacement for this selection, for a reader with access to the classical quadratic large sieve, uses two ordinary Dirichlet polynomials of length $\lfloor\sqrt X\rfloor$. At a zero, their mollified error has absolute value one. An average over the auxiliary prime shows that very few primes can have such a zero. The same argument works simultaneously for all characters of a fixed modulus, so it retains Dirichlet GRH as well as RH.

Two other steps admit shorter proofs using the estimates already established in the readable Lean source. A contour shift gives the required prime number theorem at the distant cutoff. An affine recurrence gives a better lower bound on the dimension spanned by the actual ramified forms.
\begin{center}\small\renewcommand{\arraystretch}{1.24}
\begin{tabularx}{\linewidth}{@{}p{.18\linewidth}X X@{}}
\toprule
Step & Input retained & What the replacement gives\\\midrule
Auxiliary prime & Quadratic large sieve and a character-sum bound & For $g=\beta-\Half>0$, at most $O(X^{1-g/2+\e})$ exceptional primes at fixed height. An elementary version gives $O(X^{1-g/4+\e})$.\\
Distant cutoff & The source's narrow zero-free region & A single absolutely convergent Perron integral gives $|R(x)|\ll\exp[-(\log x)^{1/27}/3]$ eventually.\\
Independent forms & Both geometric moment bounds and the ramified dimension & Missing dimension at most $20d_T/b^4$ within the original budget of translates.\\
\bottomrule
\end{tabularx}\end{center}
\textbf{Scope.} The lower bound for the covolume of long forms, the construction of short integral forms, and the geometric moment estimates still supply the arithmetic content. These substitutions give proofs on paper with explicit interfaces to the source. The published Lean checks apply to the original development. Section~\ref{sec:sources} records the extent of verification and the arithmetic that the accompanying scripts check.

\textbf{Reading routes.} The auxiliary-prime proof is in Section~\ref{sec:selector}, followed immediately by the zero contradiction. Section~\ref{sec:perron} can be read independently with the definition of the smoothed count on page~\pageref{sec:setup}. Section~\ref{sec:affine-spanning} needs only finite-dimensional linear algebra once its two moment inequalities are given.

\clearpage
\section{Where the three changes enter}\label{sec:setup}
Write $Q=18480q$ for a target character of modulus $q$, and $a=Q/2-1$. For an auxiliary prime $p\equiv a\pmod Q$, write $\eta_p(n)=(n/p)$. This class is invertible modulo $Q$ and is $7\pmod8$. For $u=1,Q$ the source uses
\begin{align}
 \lambda_{u,p}(n)&=\mathbf1_{n\equiv1\ (u)}\Lambda(n)(1-\eta_p(n)),\nonumber\\
 R_{u,p}(x)&=\sum_{n\le x}\lambda_{u,p}(n)(1/n-1/x)
                       -\frac{\log x}{\varphi(u)}-c_{u,p}.\label{eq:R}
\end{align}
Here $c_{u,p}$ is fixed by the regularized logarithmic derivative at one. A split prime contributes zero, an inert prime contributes twice its logarithm at its odd powers, and the ramified prime has the intermediate weight. This is why the auxiliary character occurs in both the arithmetic and the analytic portions.

\textbf{The near cutoff.} In the guide's notation, $m$ is the length of a chain of lifts, $a_0$ is a long derivative order (written $a$ in the source), $b$ is a weight, and $A,B$ are sets of places. The identity for the total remainder $\mathcal R$ and the volume inequality give
\begin{align*}
 \mathcal R/a_0
   &=|B|\big(R_{Q,p}(x)-R_{Q,p}(x_1)\big)+|A|R_{1,p}(x_1),\\
 m\mathcal R&\le f a_0(\log b)^5,\qquad
 f=\varphi(Q)|A|,\qquad |A|/|B|\le\log b .
\end{align*}
The inequality compares the upper lengths of short integral forms with a lower bound for the covolume of a family spanning almost the whole target. Section~\ref{sec:affine-spanning} chooses that family directly from the ramified source. It meets the same budget for denominators. Dividing the two displays gives
\[
 R_{Q,p}(x)\le \frac{\varphi(Q)(\log b)^6}{m}
             +|R_{Q,p}(x_1)|+(\log b)|R_{1,p}(x_1)|.
\]
For each sufficiently small fixed $\eta>0$ the source supplies data with $m\ge\sqrt x$, $\log b\le2x^\eta$ and $\log x_1\ge x^{\eta/12}$. Every fixed inverse logarithmic power at $x_1$ is available by Section~\ref{sec:perron}. Choosing the power large enough makes the last two terms negligible. Taking $\eta$ small, and using the source's $O(1/x)$ passage from integer to real $x$, gives
\begin{equation}
 R_{Q,p}(x)\le x^{-1/2+\e}\quad\text{eventually, for each fixed admissible } Q,p,\e>0 .
 \label{eq:arith}
\end{equation}
\textbf{The choice of the field.} An offending zero $\rho$ determines $\Half<\beta<\Re\rho$ and a height $H=|\Im\rho|$. Section~\ref{sec:selector} supplies a prime $p\equiv a\pmod Q$ such that every twist $L(s,\psi\eta_p)$ is zero-free in the box $\Re s\ge\beta$, $|\Im s|\le H$. Only after this choice is \eqref{eq:arith} invoked. Its threshold may depend on $p$ and on $\e$. Section~\ref{sec:endgame} gives the resulting contradiction. Even for zeta one takes $Q=18480$ when using these arithmetic estimates.

\clearpage
\section{Selecting the auxiliary prime with two short polynomials}\label{sec:selector}
For $Q$ a fixed positive integer, $H\ge0$ and $1/2<\beta<1$, let $\mathcal B(X)$ be the odd primes $X<p\le2X$, $p\nmid Q$, for which at least one $\psi\bmod Q$ has a zero of $L(s,\psi\eta_p)$ with $\Re s\ge\beta$ and $|\Im s|\le H$. Count each such prime once.
\begin{theorem}\label{th:selector}
For every $\e>0$,
\[
 |\mathcal B(X)|\ll_{Q,H,\beta,\e} X^{1-(\beta-1/2)/2+\e}.
\]
Consequently every fixed invertible class modulo $Q$ contains infinitely many primes at which all the twists avoid the box.
\end{theorem}
The introduction to \Source{Lib/Analytic/QuadraticTwist}{7}{the source module for quadratic twists} explicitly points out that a large sieve for real characters implies a density estimate. The calculation here identifies the version of the sieve needed, gives a first-moment proof, and supplies the exponent. In the notation $g=\beta-1/2$ the source proves $O(X^{1-g/18})$ for the admissible class. We use the specific corollary for unrestricted integer coefficients, so that prime powers and repeated prime factors are covered.

\textbf{One dyadic estimate.} Heath-Brown's Corollary~3 \cite{HB} gives, for real primitive characters of conductor at most $D$,
\[
 \sum_{\chi}\Big|\sum_{n\le N} a_n\chi(n)\Big|^2
 \ll_{e} D^e N^{1+e}(D+N)\max_{n\le N}|a_n|^2
 \qquad(e>0).
\]
The $a_n$ may be complex. Below $\|Y_p\|_2$ denotes the unnormalized $\ell^2$ norm over the primes $X<p\le 2X$, $p\nmid 2Q$. Fix a compact rectangle $K$ with left edge $\sg=1/2+\delta$, $0<\delta<\beta-1/2$, containing a neighborhood of $[\beta,1]+i[-H,H]$. For $V\le X$, $V<n\le2V$ and any deterministic $|c_n|\le\tau(n)$, we obtain, for $s\in K$,
\begin{equation}
 \bigg\|\sum_{V<n\le2V} c_n\psi(n)\eta_p(n)n^{-s}\bigg\|_2
      \ll_{e,K} X^{1/2+e} V^{-\delta}.\label{eq:block}
\end{equation}
Indeed take $D=\lceil 2X\rceil$, $N=\lfloor2V\rfloor$ and zero off the block. The coefficient bound is $\ll_{e'}V^{-\sg+e'}$ by the elementary divisor bound. Use $D+N\ll X$ and take the small sieve and divisor exponents in terms of $e$. Restricting a block to any subinterval is allowed. Summing norms on blocks above $V=U$ costs only a constant depending on $\delta$, since $\sum_{j\ge0}2^{-j\delta}<\infty$. The coefficient at $n=1$, when present, costs at most $\sqrt{2X}$.

\textbf{Ordinary truncation.} Put $\xi_p=\psi\eta_p$. Splitting the sum of $\xi_p(n)$ on any interval into classes modulo $Q$ reduces it to at most $Q$ translated sums of the Legendre character modulo $p$. Indeed $\eta_p(a+Qj)=\eta_p(Q)\eta_p(j+a Q^{-1})$ with the inverse taken modulo $p$. Pólya--Vinogradov bounds each interval sum by $O(\sqrt p\log(2p))$. This bound follows by expanding the interval in additive characters, using the size $\sqrt p$ of the nonzero Gauss sums and summing the resulting harmonic series. This treatment includes imprimitive $\psi$ and ensures that $\xi_p$ is nonprincipal. Partial summation on the convergent Dirichlet series for $\Re s>0$ gives, uniformly on $K$,
\begin{equation}
 L(s,\xi_p)-\sum_{n\le X}\xi_p(n)n^{-s}
       \ll_{Q,K} X^{1/2-\Re s}\log(2X).\label{eq:truncate}
\end{equation}
\newpage\textbf{The first moment.} Set $U=\lfloor\sqrt X\rfloor$ and form
\[
 P_p(s)=\sum_{n\le U}\xi_p(n)n^{-s},\qquad
 M_p(s)=\sum_{m\le U}\mu(m)\xi_p(m)m^{-s}.
\]
Here $M_p$ is a finite approximation to the Dirichlet inverse. The coefficient of $\xi_p(n)n^{-s}$ in $P_pM_p$ is
\begin{equation}
 b_U(n)=\sum_{\substack{m\mid n,\ m\le U\\ n/m\le U}}\mu(m).
 \label{eq:mob}
\end{equation}
It is one at $n=1$, zero for $1<n\le U$, and zero for $n>U^2$. Thus $P_pM_p-1$ is supported entirely on $U<n\le X$, and $|b_U(n)|\le\tau(n)$. Equation~\eqref{eq:block}, followed by \eqref{eq:truncate} for the tail beyond $X$, gives the three estimates
\begin{alignat}{2}
 \|M_p\|_2&\ll X^{1/2+e}, &\quad&
       \text{blocks starting at }1,\nonumber\\
 \|P_pM_p-1\|_2&\ll X^{1/2+e}U^{-\delta}, &&
       \text{blocks starting at }U,\label{eq:three}\\
 \|L(s,\xi_p)-P_p\|_2&\ll X^{1/2+e}U^{-\delta}, &&
       \text{blocks from }U\text{ to }X\text{ and the tail}.\nonumber
\end{alignat}
For the last line the tail costs $O(X^{1/2-\delta}\log(2X))$ in this norm. These estimates require neither a bound for $1/L$ nor a zero-free hypothesis.

Use $LM_p-1=(L-P_p)M_p+(P_pM_p-1)$. Cauchy--Schwarz over the varying prime, twice, gives
\begin{align}
 \sum_p|L(s,\xi_p)M_p(s)-1|
 &\le \|L-P_p\|_2\|M_p\|_2+\sqrt{2X}\,\|P_pM_p-1\|_2\nonumber\\
 &\ll X^{1+2e} U^{-\delta}\ll X^{1-\delta/2+2e}.
 \label{eq:first}
\end{align}
Using a first absolute moment at this point allows $U$ as large as $\sqrt X$. Bounding $M_p$ individually by the sum of its absolute coefficients would introduce a power of $U$ precisely where the saving is needed.

\textbf{From values to zeros.} A zero $z$ in the box belongs to $[\beta,1]+i[-H,H]$ by the Euler product for $\Re s>1$. Choose a single radius $r>0$ such that every disk $D(z,r)$ around this rectangle lies in $K$. At $z$, $|L(z,\xi_p)M_p(z)-1|=1$. Holomorphy and the disk mean-value inequality yield
\[
 \pi r^2\le\int_K |L(s,\xi_p)M_p(s)-1|\,dA(s).
\]
This compares the values at different zero locations with the same integral. Sum over bad primes for one $\psi$, integrate \eqref{eq:first}, and sum over the fixed finite set of $\psi$. Take $\sg$ sufficiently close to $\beta$ and $e$ sufficiently small. This proves the asserted bound. The prime number theorem in a fixed invertible class now supplies good primes in every sufficiently large dyadic interval. Alternatively, the source uses just the divergence of $\sum_{p\equiv a\ (Q)}(\log p)/p$ for the existence statement. A power saving makes this sum over bad primes converge. \hfill$\square$

At $\beta=3/4$ the resulting exponent is $7/8+\e$, compared with $71/72$ in the source. Section~\ref{sec:fourier} gives $15/16+\e$ in this example using only short quadratic sums and Fourier completion. These exponents count exceptional primes at a fixed height. The square-root estimate in \eqref{eq:arith} has a different origin.
\clearpage
\section{Why this choice of prime still proves RH}\label{sec:endgame}
Here is the logical check on the substitution. Assume some $L(s,\chi)$ has a zero $\rho$ with $1/2<\Re\rho<1$. After raising its modulus from $q$ to $Q=18480q$, the extra factors $1-\chi(\ell)\ell^{-s}$ are nonzero for $\Re s>0$, so the induced character still has a zero at $\rho$. Choose $\beta$ between $1/2$ and $\Re\rho$. Theorem~\ref{th:selector} supplies a prime in the required class for height $|\Im\rho|$.

Let $E_\psi$ be $L(s,\psi)$ for nonprincipal characters and the entire regularization of $(s-1)L(s,\psi)$ for the principal character. Define
\[
 N(z)=\sum_{\psi\bmod Q}\operatorname{ord}_z E_\psi,\qquad
 N_p(z)=\sum_{\psi\bmod Q}\operatorname{ord}_z E_{\psi\eta_p}.
\]
These are multiplicities. The class $1\bmod Q$ gives coefficient $+1$ for each $\psi$ by character orthogonality. The transform $\Phi(s)=\int_0^\infty R_{Q,p}(e^v)e^{(1-s)v}\,dv$ starts in a right half-plane. It continues meromorphically to $\Re s>0$, with at most simple poles, and for $z\ne 1$ its residue there is
\begin{equation}
 r(z)=\frac{N(z)-N_p(z)}{\varphi(Q)z(1-z)}.\label{eq:res}
\end{equation}
For the sign, the Dirichlet series of the weights is an average of
$L'_{\psi\eta_p}/L_{\psi\eta_p}-L'_\psi/L_\psi$, and the smoothing divides by $s(s-1)$. The coefficient of an untwisted zero is therefore $+1/[\varphi(Q)z(1-z)]$. At $1$ the centering constant cancels the would-be residue. These are the source's transform and residue calculations \cite{primeerror}.

In particular $r(\rho)\ne0$ because $N(\rho)\ge1$ and $N_p(\rho)=0$. For every real $u\in(\beta,1)$, the same prime gives $N_p(u)=0$ and $r(u)\ge0$. Above one the Euler products have no zeros. To use these signs, take $K_0>0$ large enough, by \eqref{eq:arith} with $\e=\beta-1/2$, that
\[
 d\mu(v)=\bigl(K_0e^{\beta v}-e^v R_{Q,p}(e^v)\bigr)\,dv\ge0
       \quad (v\ge0).
\]
Continuity on bounded intervals permits increasing $K_0$ once. The Laplace transform equals $K_0/(s-\beta)-\Phi(s)$ far to the right. If its real abscissa of convergence were $u>\beta$, its possible simple pole there would have residue $-r(u)\le0$. Positivity of the transform on the right forces that residue to be nonnegative, so the singularity is removable. Landau's lemma forbids this for a positive measure at its finite abscissa of convergence. Thus the integral is holomorphic on $\Re s>\beta$, forcing $r(\rho)=0$, a contradiction.

The classical boundary nonvanishing results handle $\Re s=1$ with the principal pole excepted. The functional equation then puts every zero in the open strip on its critical line, and for zeta leaves the familiar trivial zeros at the negative even integers. This is the same public target as the readable formal proof. The analytic replacement supplies the prime before the determinant estimate is applied to it. It requires no information about zeros on the critical line itself.

\clearpage
\section{A contour shift at the distant cutoff}
\label{sec:perron}

A separate shortening is available where the arithmetic argument estimates
the remainder at its distant cutoff. Because the actual smoothing vanishes
at the endpoint, its Mellin kernel decays quadratically on the initial
vertical line. This permits direct contour inversion. It also gives a
quantitative improvement over the inverse logarithmic powers stated in
\texttt{eventually\_abs\_primeRemainder\_le}.

In this section write $q$ for either modulus $1$ or $Q$ in \eqref{eq:R}, and $d$ for its fixed admissible odd prime. In particular $(d,q)=1$. Write $a_q=1/\varphi(q)$,
$\lambda(n)=\mathbf1_{n\equiv1\bmod q}\Lambda(n)(1-(n/d))$, and
$\mathcal A(s)=\sum_n\lambda(n)n^{-s}$ for $\Re s>1$.
The repository's logarithmic derivative identity gives
$\mathcal A(s)=a_qs/(s-1)+G(s)$. Put $F(s)=G(s)/s$.
The zero-free-region lemma
\texttt{eventually\_norm\_logDerivMean\_div\_le}
in \texttt{New/Reduction/Zeros/PrimeNumberTheorem.lean}
gives, for all sufficiently large $\ell$ depending on $q,d$,
\[
 F \text{ holomorphic},\qquad |F(s)|\le\ell^{26}
 \qquad (\Re s\ge1-\ell^{-26},\ |s|\le e^\ell).
\]
We claim that the centered remainder satisfies
\begin{equation}
 |R_q(x)|\ll_{q,d}\exp\bigl(-\tfrac13(\log x)^{1/27}\bigr).
 \label{eq:perron-improvement}
\end{equation}

Set $U=\log x$ and $\sigma_0=1+1/U$. On this line absolute convergence
and elementary Mellin inversion give
\[
 \sum_{n\le x}\lambda(n)(1/n-1/x)
 =\frac1{2\pi i}\int_{\sigma_0-i\infty}^{\sigma_0+i\infty}
           \frac{\mathcal A(s)x^{s-1}}{s(s-1)}\,ds .
\]
In fact the inverse transform of $y^s/[s(s-1)]$ on a line right of $1$
is $y-1$ for $y\ge1$ and zero for $0<y<1$, by closing the contour
left or right. The $y=1$ value follows by continuity and absolute
convergence. Apply this identity with $y=x/n$ and divide by $x$.
The coefficients obey $0\le\lambda(n)\le2\Lambda(n)$, so
$|\mathcal A(\sigma_0+it)|\ll U$ uniformly in $t$.
This follows either by partial summation from Chebyshev's bound or from
the simple pole of $-\zeta'/\zeta$ at $1$. Truncation at height $T\ge2$
therefore costs $O(U/T)$.

Choose $\ell=U^{1/27}$, $T=e^\ell/2$, and $w=\tfrac12\ell^{-26}$.
Shift the truncated integral to $\Re s=1-w$. The entire closed
rectangle lies strictly inside the stated region for sufficiently large
$x$. Its only pole has residue $a_q U+F(1)$, since the continued
integrand is
\[
        x^{s-1}\left(\frac{a_q}{(s-1)^2}
                       +\frac{F(s)}{s-1}\right).
\]
This is exactly the centering in $R_q(x)$. On the horizontal sides the
length is at most $1$, $|x^{s-1}|\le e$ and $|s-1|\ge T$, giving
$O(\ell^{26}/T)$ for their total. On the left side the absolute integral is
\[
 \ll e^{-wU}\left(w^{-1}+\ell^{26}\log(2+2T/w)\right),
\]
by integrating $(w^2+t^2)^{-1}$ and $(w^2+t^2)^{-1/2}$ for
$|t|\le T$. Here $wU=\ell/2$ and $U=\ell^{27}$, so all three errors
are $O(e^{-\ell/3})$. This proves \eqref{eq:perron-improvement}.
It supplies every fixed inverse logarithmic power needed in the
two-cutoff comparison and dispenses with the gamma-kernel localization
in this application. The narrow classical zero-free region still
provides the shift. The square-root gain at the near cutoff still comes
from the theta chain.

\clearpage
\section{Choosing translates directly on the ramified source}
\label{sec:affine-spanning}

The volume comparison needs a small family of translates whose images span
almost all of the target. The source first chooses translates of the entire
space of Poincar\'e series. It then restricts each map to the ramified
source and subtracts a possible dimension loss for every translate. One
can restrict at each choice instead. This gives the following stronger
version of the spanning conclusion. Here $b$ is the weight parameter,
$d$ is the target dimension, and $r$ is the ramified source dimension.
For $b\ge 8$ there is a set $J$ of translates satisfying
\begin{equation}
 |J|r\le 5d\log(1+2b^2),\qquad
 \operatorname{codim}\sum_{s\in J}
  \operatorname{im}(\text{normalized Poincar\'e map at }s)
 \le \frac{20d}{b^4}.
 \label{eq:affine-spanning-bound}
\end{equation}
The volume argument asks for codimension at most $4d/b^2$ with exactly
this translate budget. The proof below uses the existing geometric moment
estimates at their original precision.

\subsection*{An affine recurrence for the missing dimension}

Let $S$ and $V$ have positive dimensions $h$ and $d$. Let $P_s:S\to V$
be a nonempty finite family of linear maps and let $S_0\subset S$ have
dimension $r\ge (1-\tau)h$, where $0\le\tau\le1$. Suppose that for
every $F\subset V$ of codimension $c$ there is an $s$ with
\begin{equation}
 \dim P_s^{-1}(F)\le h\left(1-\frac{c}{2d}+\omega\right).
 \label{eq:affine-spread}
\end{equation}
Assume $0<h\le 2d$ and $\omega\ge0$. For each integer $u\ge 0$
one can choose at most $u$ indices whose restricted images have
codimension at most
\begin{equation}
  d\left[\kappa+(1-\kappa)\left(1-\frac{h}{2d}\right)^u\right],
  \qquad \kappa=2\omega+2\tau.
  \label{eq:affine-spread-recurrence}
\end{equation}
In particular, the missing fraction is at most
$\kappa+\exp(-uh/(2d))$.

To prove this, begin with $F=0$. If $F$ is the current span, use
\eqref{eq:affine-spread} and compose $P_s|_{S_0}$ with the quotient
$V\to V/F$. Its kernel is $S_0\cap P_s^{-1}(F)$. The dimension it adds is
therefore at least
\[
 r-\dim P_s^{-1}(F)\ \ge\ \frac{h}{2}\left(\frac{c}{d}
                                      -2\omega-2\tau\right).
\]
Writing $q_j$ for the missing fraction after $j$ choices and
$t=h/(2d)$, we have $q_{j+1}\le(1-t)q_j+t\kappa$ and $q_0=1$.
Iteration gives \eqref{eq:affine-spread-recurrence}, including when the
one-step lower bound is negative near full coverage. Finally
$(1-t)^u\le e^{-tu}$. Repeated indices contribute only one element to
the resulting set. The loss $\tau$ enters the limiting missing fraction.
It is absorbed into the fixed point of the recurrence.

\subsection*{Checking the hypotheses against the source}

Write $P$ for the scaled full-source Poincar\'e map, $e$ for orthogonal
projection onto the target, and $s$ for the unitary action of a member
of the finite group. The target is invariant and $e$ commutes with the
action. The geometric estimates in
\texttt{SpanningTheorem.lean}, at
\Source{New/Spanning/SecondMoment/SpanningTheorem}{257}{the projected near-isometry} and
\Source{New/Spanning/SecondMoment/SpanningTheorem}{288}{the normalized second moment},
give, in Hilbert--Schmidt norm,
\[
 \|eP\operatorname{pr}_U\|_{\rm HS}^2\ge\dim U-5b^{-4}h
 \quad(U\subset S),\qquad
 d\,\underset{s}{\operatorname{avg}}\,\|P^*sP\|_{\rm HS}^2
 \le(1+9b^{-4})h^2 .
\]
In particular $U=\ker(eP)$ gives $d\ge(1-5b^{-4})h$. Thus $h\le 2d$
for $b\ge8$. Write $P_s=esP$.
For a subspace $F$ of relative dimension $a=1-c/d$, put
$K=|G|^{-1}\sum_s sPP^*s^*$. Cycling the trace gives
\[
 \|K\|_{\rm HS}^2=\underset{s}{\operatorname{avg}}\,\|P^*sP\|_{\rm HS}^2,
 \qquad
 \underset{s}{\operatorname{avg}}\,\|\operatorname{pr}_F sP\|_{\rm HS}^2
 =\operatorname{tr}(\operatorname{pr}_F K)
 \le h\sqrt{a(1+9b^{-4})}.
\]
Choose $s$ whose projected energy is at most this average. On
$U=P_s^{-1}(F)$ one has $esP=\operatorname{pr}_F sP$, and unitarity
and $es=se$ give
\[
 \dim U-5b^{-4}h
 \le \|eP\operatorname{pr}_U\|_{\rm HS}^2
 \le \|\operatorname{pr}_F sP\|_{\rm HS}^2.
\]
Now $2\sqrt{a(1+9b^{-4})}\le a+1+9b^{-4}$ proves
\eqref{eq:affine-spread} with $\omega=(5+9/2)b^{-4}$.
This is the trace argument of
\Source{New/Spanning/SecondMoment/FewTranslates}{320}{the one-translate lemma}
with the two errors kept separate.

The ramified codimension estimate used in
\Source{New/Reduction/ShortAndEnough}{410}{the published choice of translates} is
$h-r\le b^{-8}h$. The source-coordinate change transports $S_0$;
the nonzero normalization factors preserve its image. Apply the
recurrence with $\tau=b^{-8}$ and
\[
        u=\left\lceil\frac{2d}{h}\log(b^5)\right\rceil .
\]
The codimension fraction is at most $19b^{-4}+2b^{-8}+b^{-5}\le20b^{-4}$. For the count, use $r\le h$, $h\le2d$ and $b\ge8$ to obtain
\[
 |J|r\le 10d\log b+h\le 10d\log b+2d
 \le 5d\log(2b^2)\le5d\log(1+2b^2).
\]
This proves \eqref{eq:affine-spanning-bound}. Since
$20/b^4\le4/b^2$ in this range, Theorem B supplies its published
interface to the determinant comparison. A minimal formal substitution can use the existing one-translate lemma verbatim: taking $\varepsilon=9b^{-4}$ gives $\omega=3\varepsilon/2$ and codimension at most $29d/b^4$, also sufficient. The $20$ uses the separate errors above. The new step is the recurrence on
the restricted image and its iteration. The stronger estimate does
not change the exponent in the prime sum, which also depends on the
length and cost of the theta chain.

\subsection*{Why the number of translates contains a logarithm}

There is a finite example with ideal moments in which the logarithm
has the correct order. Give $\mathbb C^{\binom Nk}$ coordinates indexed
by the $k$-subsets of an $N$-element set. The symmetric group permutes
the coordinates. Let $P$ include the coordinate subspace spanned by
the $k$-subsets containing a specified point. Here $d=\binom Nk$,
$h=\binom{N-1}{k-1}$, and $h/d=k/N$. The translated image picks a
point and spans all coordinates containing it. Each coordinate occurs
in $k$ of the $N$ images, so the average of their orthogonal projectors
is $(h/d)I$. In particular the first moment is exact and the mean of
$\|P^*sP\|_{\rm HS}^2$ is $h^2/d$.

Choosing $j$ distinct points leaves exactly
$\binom{N-j}{k}/\binom Nk$ of the coordinates uncovered. For fixed $L>0$,
take $k\to\infty$, $k/N\to0$, and $j=\lfloor LN/k\rfloor$. The logarithm
of this fraction is
\[
 \sum_{a=0}^{k-1}\log\left(1-\frac{j}{N-a}\right)=-L+o(1).
\]
Indeed the sum of the fractions tends to $L$ and the sum of their
squares tends to zero. Thus even exact moments can require order
$(d/h)\log(1/\theta)$ translates to reach a prescribed missing
fraction $\theta$. This example concerns the information in the two
moment bounds. The particular Poincar\'e maps could have additional
structure.


\clearpage
\section{Source interfaces and verification}\label{sec:sources}
The 79-page attached guide describes the readable \texttt{grh/} package. Its sections 9, 16, 17 and 18 describe the moments, the two cutoffs, Landau's argument and the twist selection. The source below is pinned to the guide's ProveIt commit \texttt{d8415e76}. Its readable directory, the one at the inspected main commit \texttt{8dd060b0}, and the independently checked upstream revision \texttt{72fead10} all have Git tree \texttt{2e133df63d7f4da6b3059c4195324f8e9187ba73}. The file \texttt{source-anchors.tsv} gives immutable links and Git blob hashes.
\begin{center}\small\renewcommand{\arraystretch}{1.26}
\begin{tabularx}{\linewidth}{@{}p{.21\linewidth}X X@{}}
\toprule
Proof step & Existing inputs & Replacement boundary\\\midrule
Choice of a good prime &
\Source{Lib/Analytic/QuadraticTwist}{206}{Three-lines zero detector}, \Source{Lib/Analytic/QuadraticTwist}{536}{theta-series mean square}, \Source{Lib/Analytic/QuadraticTwist}{732}{exceptional count} &
\Source{Lib/Analytic/QuadraticTwist}{796}{Existence of a simultaneous good twist}, consumed by \Source{New/Reduction/Zeros/Criterion}{52}{the one-sided criterion}.\\
Distant cutoff &
\Source{New/Reduction/Zeros/PrimeError}{263}{Prime Dirichlet series} and \Source{New/Reduction/Zeros/PrimeNumberTheorem}{178}{the narrow zero-free estimate} &
\Source{New/Reduction/Zeros/PrimeNumberTheorem}{217}{Every fixed inverse logarithmic power}, used for $q=1,Q$ in \Source{New/Reduction/OneSidedBound}{95}{the far-cutoff bound}.\\
Ramified span &
\Source{New/Spanning/SecondMoment/FewTranslates}{320}{One-translate preimage estimate} and \Source{New/UpperBound/Size/ShortColumns}{944}{ramified dimension} &
\Source{New/Reduction/ShortAndEnough}{410}{Choice of translates} and \Source{New/Reduction/ShortAndEnough}{448}{Theorem B}, used by \Source{New/Reduction/MainInequality}{87}{the main inequality}.\\
\bottomrule
\end{tabularx}\end{center}
\textbf{Which substitution to make first.} For a human proof, the two-polynomial argument is the shortest path from a standard reference to the exact good-twist statement. The original source already identifies the quadratic large sieve as an alternative and cites Jutila and Heath-Brown. The worked substitution above and its exponent make that suggestion explicit. If avoiding that external theorem matters, Appendix~\ref{sec:fourier} proves a sufficient mean estimate by completing smooth sums. The fixed-height bounds here make no claim to priority as zero-density theorems.

For a Lean refactor, the ramified-span recurrence is the smallest self-contained target. It reuses the one-translate estimate, changes a finite-dimensional argument, and improves its conclusion. The Perron proof requires a formal contour inversion and estimates for the three sides. The large-sieve proof would also require a formal version of the real-character inequality or its proof. A shorter mathematical argument can therefore have a larger formal dependency. The published kernel checks by Sutherland cover the original five statements \cite{checks}. Here the new arguments were checked on paper, the source interfaces were inspected, and finite identities were checked by the included scripts. There has been no new Lean replay of any substitution.

{\small
\begin{thebibliography}{9}
\bibitem{HB} D. R. Heath-Brown.
\emph{A mean value estimate for real character sums}.
Acta Arithmetica 72 (1995), 235--275, especially Corollary 3, p.~238.
\href{https://doi.org/10.4064/aa-72-3-235-275}{Journal article}.
\bibitem{source} Levent Alpöge and Claude. Readable Lean development for RH and Dirichlet GRH, revision \texttt{72fead10}, 2026.
\href{https://github.com/lsernau/grh-lean/tree/72fead102f663d5c7a407964a13dc0cf9d977207/grh}{Readable package}.
\Source{Lib/Analytic/QuadraticTwist}{7}{Discussion of the large-sieve alternative}.
\Source{New/Reduction/OneSidedBound}{239}{Half-plane theorem}.
\Source{Main}{5}{Public zeta theorem}.
\bibitem{primeerror} Readable development, \emph{Prime error and the one-sided criterion}.
\Source{New/Reduction/Zeros/PrimeError}{293}{Transform}.
\Source{New/Reduction/Zeros/PrimeError}{414}{Residues}.
\Source{Lib/Analytic/Landau}{217}{Landau's lemma in the required form}.
\bibitem{checks} Andrew V. Sutherland. \emph{Verification files for the shorter Lean proof of RH and of GRH for Dirichlet $L$-functions}, 27 November 2026.
\href{https://math.mit.edu/~drew/shortgrh/}{Verification record and logs}.
\bibitem{origin} Levent Alpöge. \emph{Announcement and remarks on the proof}, 2026.
\href{https://alpo.ge/remark.pdf}{The chain-length idea and its Stepanov motivation}.
\end{thebibliography}
}

\clearpage
\appendix
% Standalone mathematical insert; no macro prerequisites beyond amsmath, amssymb.
\section{Selecting the twist by elementary Fourier completion}\label{sec:fourier}
One can also prove the fixed-height twist-selection lemma using Fourier completion. A smooth partition of the Dirichlet series reduces the required character-sum estimate to P\'olya--Vinogradov. The first-moment and disk arguments from Section~\ref{sec:selector} apply again. Throughout this section the modulus $Q$, the rectangle and its slightly larger neighborhood are fixed. Let $1/2<\sigma_0=1/2+\delta<1$ be a lower bound for the real parts in the neighborhood. Let $X$ be large and let $p$ run over primes $X<p\le 2X$. Fix one primitive character $\psi$ of conductor $c\mid Q$, including $c=1$. For $p\nmid c$, the product $\xi_p=\psi(\cdot)(\cdot/p)$ is primitive of conductor $cp$ by the Chinese remainder theorem. Passing from this primitive character to the character induced modulo $Qp$ only multiplies its $L$-function by factors $1-\xi_p(\ell)\ell^{-s}$, which have no zeros in $\Re s>0$.

\paragraph{A short mean square.} For arbitrary complex $a_n$ supported on $1\le n\le V$, put $A=\max_n|a_n|$. For every $e>0$,
\begin{equation}
 \left(\sum_{X<p\le2X}\left|\sum_{n\le V}a_n(n/p)\right|^2\right)^{1/2}
 \ll_e (XV)^e\left(\sqrt{XV}+V^{3/2}\right)A.\tag{C1}
\end{equation}
To see this, extend the family on the left to all positive odd denominators up to $2X$ using the Jacobi symbol. Expand the square. When $n_1n_2$ is a square the inner sum is at most $2X$, and there are $O(V\log(2V))$ such pairs, since $n_j=kr_j^2$ with common squarefree $k$. For a positive nonsquare $n_1n_2$, quadratic reciprocity expresses the Jacobi symbol on odd denominators as a nonprincipal Dirichlet character of modulus at most $8n_1n_2$, zero on even denominators. This also treats twice a square, through the nontrivial character modulo eight. Removing imprimitive factors by inclusion--exclusion and applying P\'olya--Vinogradov bounds it by $O_e(V^{1+e})$. There are at most $V^2$ pairs. Renaming $e$ gives (C1). This calculation also appears in the introduction to Heath-Brown's 1995 paper, prior to its main large-sieve theorem.

\paragraph{Completion of the long blocks.} Let $w$ be a smooth function supported on $[1,4]$ and $Y\ge1$. Uniformly for $s=\sigma+it$ in the chosen bounded neighborhood,
\begin{equation}
 \left(\sum_{X<p\le2X}\left|\sum_{n\ge1}\xi_p(n)n^{-s}w(n/Y)\right|^2\right)^{1/2}
 \ll_{Q,w,e} X^{1/2+e}Y^{1/2-\sigma}.\tag{C2}
\end{equation}
For $Y\le\sqrt X$ this is (C1), with the fixed character and the smooth weight in the coefficients. For $Y\ge\sqrt X$ apply Poisson summation in the residue classes modulo $cp$. Set $K_s(u)=\int_0^\infty v^{-s}w(v)e^{-2\pi iuv}\,dv$. The primitive Gauss-sum formula gives the exact block
\[
 \frac{Y^{1-s}\tau(\xi_p)}{cp}
 \sum_{m\ne0}\overline{\xi_p}(m)K_s\!\left(\frac{mY}{cp}\right),
 \qquad |\tau(\xi_p)|=\sqrt{cp}.
\]
The zero frequency has vanished. For each sign of $m$ let $S_p(t)=\sum_{1\le m\le t}\overline\psi(m)(m/p)$. The parity factor in the negative part has absolute value one. Put $Z=cX/Y$. The derivative of $K_s(\pm tY/(cp))$ is bounded by $C_A Z^{-1}(1+t/(2Z))^{-A}$, uniformly in $p$ and $s$, by integration by parts in the defining Fourier integral. Partial summation, Minkowski's inequality and (C1) bound the $\ell^2$ norm of each frequency sum by
\[
 C_{A,e}X^e\int_1^\infty
 \frac{t^e}{Z}(1+t/(2Z))^{-A}
 \bigl(\sqrt{X}t^{1/2}+t^{3/2}\bigr)\,dt.
\]
For $1\le Z\le c\sqrt X$ this is $O_{c,e}(X^{O(e)}\sqrt{XZ})$, by $t=Zu$. For $0<Z<1$ use $(1+t/(2Z))^{-A}\le(2Z/t)^A$ with $A\ge5+e$, obtaining the same bound. Multiply by $Y^{1-\sigma}/\sqrt{cp}$, use $p\asymp X$, and rename $e$ to obtain (C2). The decay in (C2) for $\sigma>1/2$ is summable over increasing dyadic $Y$.

\paragraph{Mollification.} Take a smooth cutoff $\phi$ equal to one on $[0,1]$ and zero on $[2,\infty)$. Let $U=M=\lfloor X^{1/4}/4\rfloor$, and put
\[
 P_p(s)=\sum_n\xi_p(n)n^{-s}\phi(n/U),\qquad
 D_p(s)=\sum_{l\le M}\mu(l)\xi_p(l)l^{-s}.
\]
Telescope the smooth cutoffs at $U,2U,4U,\ldots$. Consecutive differences use $w(v)=\phi(v/2)-\phi(v)$ supported on $[1,4]$. For fixed $p$ the cutoffs converge to $L(s,\xi_p)$ in $\sigma>0$ by the bounded partial sums of a nonprincipal character and partial summation. Applying (C2) and summing geometrically gives
\[
 \|L(s,\xi_p)-P_p(s)\|_{\ell_p^2}\ll X^{1/2+e}U^{-\delta}.
\]
The short mean square (C1) bounds $\|D_p\|_{\ell_p^2}\ll X^{1/2+e}$. Indeed, on each dyadic block of length comparable to $V\le M$ the coefficients after extracting the symbol $(n/p)$ have modulus at most $V^{-\sigma_0}$, up to a fixed constant. Equation (C1) gives $O(X^{1/2+O(e)}V^{-\delta})$, summable over $V=1,2,4,\ldots$. Meanwhile $P_pD_p-1$ has support $M<n\le 2UM$ with coefficients of $\xi_p(n)n^{-s}$ bounded by $\|\phi\|_\infty\tau(n)$. Vanishing for $1<n\le M$ follows from $M\le U$ and $\sum_{l\mid n}\mu(l)=0$. Since $2UM\le\sqrt X$, apply (C1) on dyadic blocks $V<n\le 2V$ and sum the norms, using $\tau(n)\ll_e n^e$. This gives
\[
 \|P_pD_p-1\|_{\ell_p^2}\ll X^{1/2+O(e)}M^{-\delta}.
\]
Use the first absolute moment for the final estimate. Cauchy--Schwarz, first on the product and then on the short polynomial, gives
\begin{align*}
 \sum_{X<p\le2X}|L(s,\xi_p)D_p(s)-1|
 &\le \|L(s,\xi_p)-P_p(s)\|_{\ell_p^2}\|D_p(s)\|_{\ell_p^2}
       +\sqrt{2X}\|P_pD_p-1\|_{\ell_p^2}\\
 &\ll X^{1+O(e)}(U^{-\delta}+M^{-\delta})
 \ll X^{1-\delta/4+O(e)} .
\end{align*}
All bounds are uniform on the neighborhood. If $L$ vanishes at a point of the original compact rectangle, $|LD_p-1|=1$ there. The disk mean-value inequality for the absolute value of a holomorphic function gives a uniform positive lower bound for the integral of $|LD_p-1|$ on the larger neighborhood. Integrating the last display bounds the exceptional primes by $O(X^{1-\delta/4+O(e)})$. For a rectangle with left edge $\beta>1/2$, choose the disk margin and the losses $e$ sufficiently small in terms of a prescribed $\varepsilon>0$. The resulting count is
\[
 O_{Q,H,\beta,\varepsilon}
 \bigl(X^{1-(\beta-1/2)/4+\varepsilon}\bigr).
\]
A finite union covers all characters modulo $Q$, and the prime number theorem for the fixed admissible progression supplies a good prime.


\end{document}
